\section{关键实现细节}

\subsection{守恒量到原始量的转换}

\begin{align}
\rho &= q_1,\\
u &= \frac{q_2}{q_1},\\
p &= (\gamma-1)
\left(
q_3-\frac{q_2^2}{2q_1}
\right).
\end{align}

说明为何 $q_1>0$ 且 $p>0$ 是后续计算声速的必要条件。

\subsection{数值异常检测}

建议在最终程序中记录并检查：

\begin{itemize}[leftmargin=2em]
    \item $\rho\le 0$ 或 $p\le 0$；
    \item \texttt{NaN} 或 \texttt{Inf}；
    \item CFL 超出稳定范围；
    \item 时间步过小或达到最大迭代次数；
    \item 波到达计算域边界。
\end{itemize}

\subsection{快照与重启}

快照文件同时保存原始量和守恒量：

\begin{equation}
(x,\rho,u,p,E,q_1,q_2,q_3).
\end{equation}

说明快照用于绘图、动画和中间状态分析。若要实现严格 restart，
还应保存当前时间、步数、网格参数和气体参数，并增加读取函数。

\subsection{程序当前限制}

在此如实记录最终提交版本仍存在的限制，例如：

\begin{itemize}[leftmargin=2em]
    \item 人工粘性公式的经验性；
    \item MacCormack 对强间断的振荡；
    \item 边界条件较简单；
    \item 尚未实现高分辨率限制器；
    \item 尚未实现读取快照继续计算。
\end{itemize}
